\documentclass[11pt, a4paper]{article}

% ----------------------------------------------------------------------
% PAKETE & EINSTELLUNGEN
% ----------------------------------------------------------------------
\usepackage[utf8]{inputenc}
\usepackage[T1]{fontenc}
\usepackage[ngerman]{babel}
\usepackage{amsmath, amssymb, amsfonts, amsthm}
\usepackage{listings}
\usepackage{xcolor}
\usepackage{hyperref}
\usepackage{geometry}
\usepackage{fancyhdr}

\geometry{left=2cm, right=2cm, top=2cm, bottom=2.5cm}
\setlength{\parindent}{0pt}
\setlength{\parskip}{0.6em}

% Code-Listing Stil
\lstset{
    language=Python,
    basicstyle=\ttfamily\small,
    keywordstyle=\color{blue},
    commentstyle=\color{gray},
    stringstyle=\color{red},
    numbers=left,
    numberstyle=\tiny\color{gray},
    breaklines=true,
    frame=single,
    backgroundcolor=\color{gray!5},
    showstringspaces=false,
}

% Theorem-Umgebungen
\newtheorem{theorem}{Theorem}
\newtheorem{definition}{Definition}
\newtheorem{remark}{Bemerkung}

% Titel-Metadaten
\title{\textbf{Quaternionen-Zerlegung}\\[0.5em]
\large Zerlegung natürlicher Zahlen in Prim-Quaternionen im Hurwitz-Quaternionen-Ring}

\author{SageMath 10.8}
\date{\today}

\begin{document}

\maketitle

\begin{abstract}
\noindent
Dieser Artikel beschreibt die Zerlegung natürlicher Zahlen in ein Produkt von Prim-Quaternionen im Hurwitz-Quaternionen-Ring $\mathcal{H}$. Ausgehend vom Satz von Lagrange (jede natürliche Zahl ist Summe von vier Quadraten) wird ein Algorithmus vorgestellt, der zu gegebenem $n \in \mathbb{N}$ eine Faktorisierung in Prim-Quaternionen $\pi_1, \ldots, \pi_k$ mit $N(\pi_1) \cdots N(\pi_k) = n$ liefert. Die Implementierung erfolgt in SageMath.
\end{abstract}

% ----------------------------------------------------------------------
\section{Import und Hurwitz-Quaternionen-Ring}
% ----------------------------------------------------------------------

Der Hurwitz-Quaternionen-Ring $\mathcal{H}$ ist die Quaternionen-Algebra über $\mathbb{Q}$ mit den Parametern $(-1,-1)$, d.\,h.\, $i^2 = j^2 = k^2 = ijk = -1$ (Hamilton). In SageMath wird er definiert durch:

\begin{lstlisting}
from sage.all import *
from sage.arith.misc import four_squares

# Definition des Hurwitz-Quaternionen-Rings uber Q
H = QuaternionAlgebra(QQ, -1, -1)
H
\end{lstlisting}

Ausgabe: \texttt{Quaternion Algebra (-1, -1) with base ring Rational Field}

% ----------------------------------------------------------------------
\section{Funktion: Zerlegung in Prim-Quaternionen}
% ----------------------------------------------------------------------

Die \textbf{Quaternionen-Norm} eines Elements $q = a + bi + cj + dk$ ist
\[
N(q) = a^2 + b^2 + c^2 + d^2.
\]
Nach dem \textbf{Satz von Lagrange} ist jede natürliche Zahl $n$ Summe von vier Quadraten: $n = a^2 + b^2 + c^2 + d^2$. Somit existiert stets ein Quaternion $q \in \mathcal{H}$ mit $N(q) = n$.

Die Funktion \texttt{quaternions\_zerlegung(n)} zerlegt $n$ in eine Kette von Prim-Quaternionen, indem für jeden Primfaktor $p$ von $n$ ein Quaternion $\pi$ mit $N(\pi) = p$ gefunden wird (ebenfalls über \texttt{four\_squares}).

\begin{lstlisting}
def quaternions_zerlegung(n):
    """
    Zerlegt eine naturliche Zahl n in ein Produkt von Prim-Quaternionen.
    Basiert auf dem Hurwitz-Quaternions-Ring (H).
    """
    # Wir suchen ein Quaternion q, dessen Norm n ist
    # (Satz von Lagrange: Jede naturliche Zahl ist Summe von 4 Quadraten)
    coords = four_squares(n)
    q = H(coords)
    
    print(f"### #Energiedoku: Analyse von n = {n}")
    print(f"Basis-Quaternion q: {q}")
    print(f"Norm N(q): {q.reduced_norm()}")
    print("-" * 30)
    
    # Zerlegung der Norm n in Primfaktoren p_i
    # Fur jedes p_i ein zugehoriges Prim-Quaternion pi_i
    faktoren_n = factor(n)
    prim_quaternionen_kette = []
    
    for p, expo in faktoren_n:
        p_coords = four_squares(p)
        pi = H(p_coords)
        for _ in range(expo):
            prim_quaternionen_kette.append(pi)
            
    return prim_quaternionen_kette
\end{lstlisting}

% ----------------------------------------------------------------------
\section{Beispiel}
% ----------------------------------------------------------------------

Für $n = 110567868768768768767$ erhält man:

\begin{lstlisting}
n_test = 110567868768768768767
kette = quaternions_zerlegung(n_test)

print("Zerlegung in Prim-Quadrupel (Familien-Komponenten):")
for i, pi in enumerate(kette):
    print(f"Faktor {i+1} (Prim-Quaternion): {pi} | Norm: {pi.reduced_norm()}")
\end{lstlisting}

\textbf{Ausgabe:}
\begin{verbatim}
### #Energiedoku: Analyse von n = 110567868768768768767
Basis-Quaternion q: 66 + 1511*i + 92589*j + 10515125713*k
Norm N(q): 110567868768768768767
------------------------------
Zerlegung in Prim-Quadrupel (Familien-Komponenten):
Faktor 1 (Prim-Quaternion): 3 + 5*i + 177*j + 16188*k | Norm: 262082707
Faktor 2 (Prim-Quaternion): 1 + 20*i + 402*j + 649524*k | Norm: 421881588581
\end{verbatim}

% ----------------------------------------------------------------------
\section{Verifikation}
% ----------------------------------------------------------------------

Das Produkt der Normen der Prim-Quaternionen muss wieder $n$ ergeben:

\begin{lstlisting}
check_norm = prod(p.reduced_norm() for p in kette)
print("-" * 30)
print(f"Verifikation der Gesamt-Energie (Produkt der Normen): {check_norm}")
print(f"Erwartet: {n_test} -> Korrekt: {check_norm == n_test}")
\end{lstlisting}

\textbf{Ausgabe:}
\begin{verbatim}
------------------------------
Verifikation der Gesamt-Energie (Produkt der Normen): 110567868768768768767
Erwartet: 110567868768768768767 -> Korrekt: True
\end{verbatim}

% ----------------------------------------------------------------------
\section*{Zusammenfassung}
% ----------------------------------------------------------------------

Die Zerlegung natürlicher Zahlen in Prim-Quaternionen im Hurwitz-Ring nutzt den Satz von Lagrange und die Faktorisierung in $\mathbb{Z}$. Für jeden Primfaktor $p$ von $n$ wird ein Quaternion $\pi$ mit $N(\pi) = p$ gefunden; das Produkt der Normen aller Prim-Quaternionen ergibt $n$ und bestätigt die Korrektheit der Zerlegung.

\end{document}
